%% kapa: exam=1 round=25 year=2014 period=1 form=A subject=civil_law subject_ko=민법 seq=27 number=27
%% source: Q-Net 감정평가사 기출문제 (cst00309) artlSeq=5101207 file=extracted/1008/5101207/2014년 제25회 감정평가사 1차시험 기출문제/1교시(A형).hwp
%% conversion: kordoc 4.24.51 -> tools/gen1.py (auto); review=proofread-sonnet; flags=none
\begin{question}{27}
소유권에 관한 설명으로 옳지 않은 것은? (다툼이 있으면 판례에 의함)
\begin{choices}
  \item 소유자는 법률의 범위 내에서 소유물을 사용, 수익, 처분할 권리가 있다.
  \item 지적공부에 등록된 토지의 소유권 범위는 특별한 사정이 없는 한, 지적공부상 경계에 의하여 확정된다.
  \item 구분소유는 구분행위가 행해지고 건물이 완성되었더라도 집합건축물대장의 등록이나 구분건물의 표시에 관한 등기가 완료된 때에 성립한다.
  \item 구분소유권의 객체로서 적합한 물리적 요건을 갖추지 못한 건물의 일부가 구분소유권의 목적으로 등기되었더라도 구분소유권은 인정되지 않는다.
  \item 유실물은 법률에 정한 바에 의하여 공고한 후 6개월 내에 그 소유자가 권리를 주장하지 아니하면 다른 사정이 없는 한, 습득자가 그 소유권을 취득한다.
\end{choices}
\end{question}
\answer{3}
